\chapter{Multi Round-Trip Requests}
\label{ch:mrtr}

\epigraph{``What we've got here is failure to communicate.''}{The Captain,
\emph{Cool Hand Luke} (1967)}

The Captain says it right before demonstrating that he has no interest in
communicating at all. Which is roughly what a 2025-era MCP server did when it
needed something from you: it stopped, held the connection hostage, and sent its
own request back down the pipe.

That mechanism is gone. What replaced it is one of the cleanest pieces of design
in the whole revision, and it costs more than people expect. Both halves of that
sentence matter, so this chapter does the mechanism first and the price second.

\section{The problem}

A server is halfway through a tool call and needs something only the client can
supply. A username. An approval. A confirmation from a human who has authority
the server does not.

The old answer was control-flow inversion. The server sent a JSON-RPC
\emph{request} back to the client, mid-call, on the same connection, and blocked
until the answer arrived.

\bookfig[0.90]{mrtr-sequence}{Above, the server holds a request open and owns
state for its duration. Below, the request finishes, the state travels with the
retry, and a different replica can answer.}{mrtr-sequence}

That worked, and it had a fatal property for 2026: the connection meant
something. The server held state for the life of the request. Which meant no
load balancing without stickiness, no serverless, and no surviving a dropped
connection.

Statelessness and server-initiated requests cannot coexist. One of them had to
go.

\section{The MRTR pattern}

Four steps, and the third one is the clever bit.

\begin{enumerate}[leftmargin=1.6em, itemsep=3pt]
\item Client sends a request.
\item Server realises it needs more, and \textbf{answers}. Not with the result,
with \texttt{resultType: "input\_required"}, carrying the questions and an opaque
\texttt{requestState}. The original request is now complete and forgotten.
\item Client gathers what was asked for, then \textbf{re-sends the original
request}, with a \emph{new} id, carrying \texttt{inputResponses} and echoing
\texttt{requestState} unchanged.
\item Server has enough, and returns the real result.
\end{enumerate}

The server holds nothing between steps 2 and 3. Whatever it needs to remember,
it seals into \texttt{requestState} and hands to the client to carry back.

\subsection{On the wire}

The interim result:

\begin{wire}
{
  "jsonrpc": "2.0",
  "id": 1,
  "result": {
    "resultType": "input_required",
    "inputRequests": {
      "approval": {
        "method": "elicitation/create",
        "params": {
          "mode": "form",
          "message": "Exposure of $9,000,000 on ACC-1042 exceeds the
                      $5,000,000 threshold. Who is approving?",
          "requestedSchema": {
            "type": "object",
            "properties": {
              "approver": { "type": "string", "title": "Approver",
                            "minLength": 2 },
              "rationale": { "type": "string", "title": "Rationale",
                             "maxLength": 400 }
            },
            "required": ["approver"]
          }
        }
      }
    },
    "requestState": "v1.eyJkIjp7ImFjY291bnRJZCI6...?=.mXK9Fz..."
  }
}
\end{wire}

And the retry:

\begin{wire}
{
  "jsonrpc": "2.0",
  "id": 2,
  "method": "tools/call",
  "params": {
    "name": "assess_account_risk",
    "arguments": { "accountId": "ACC-1042", "exposureUsd": 9000000 },
    "inputResponses": {
      "approval": {
        "action": "accept",
        "content": { "approver": "j.okonjo", "rationale": "board approved" }
      }
    },
    "requestState": "v1.eyJkIjp7ImFjY291bnRJZCI6...?=.mXK9Fz..."
  }
}
\end{wire}

Three things to notice, all of them enforced:

\textbf{The id changed.} 1 to 2. These are independent requests, not a
continuation, and the specification requires the ids to differ. Meridian tests
it:

\begin{py}
def test_retry_uses_a_different_id(self):
    """The retry is an independent request, so it must not reuse the id."""
    ...
    result = client.call_tool("ask", {})
    self.assertEqual(result["content"][0]["text"], "hello ada")
    self.assertEqual(len(ids), 2)
    self.assertNotEqual(ids[0], ids[1])
\end{py}

\textbf{The arguments are repeated in full.} The server does not remember them.

\textbf{\texttt{requestState} is echoed byte for byte.} The client must not
inspect it, parse it, or modify it. If there was no \texttt{requestState}, the
client must not invent one.

\subsection{The three requests a server may make}

\texttt{inputRequests} is a map from server-chosen keys to request objects, and
exactly three kinds are permitted:

\begin{table}[htbp]
\centering
\small
\begin{tabularx}{\textwidth}{@{}lXl@{}}
\toprule
\thead{Method} & \thead{Asks for} & \thead{Status} \\
\midrule
\texttt{elicitation/create}   & Input from the \emph{user} & Active \\
\texttt{sampling/createMessage} & A completion from the \emph{model} & Deprecated \\
\texttt{roots/list}           & The client's directory roots & Deprecated \\
\bottomrule
\end{tabularx}
\caption{Two of the three are on their way out. Elicitation is the one to build
on.}
\label{tab:input-requests}
\end{table}

A map rather than a list, because a server can ask several things at once and
the keys correlate answers to questions. Asking for everything in one interim
result costs one round trip; asking one at a time costs one each.

And a hard rule: \textbf{never send a request the client did not declare support
for.} Meridian checks and degrades:

\begin{py}
def _ask_for_approval(self, ctx, account_id, exposure):
    if not ctx.capabilities.supports_elicitation("form"):
        # Never send a request shape the client did not declare. Degrade
        # to a plain error the model can relay to the user instead.
        return tool_error(
            f"Exposure of ${exposure:,.0f} exceeds the "
            f"${APPROVAL_THRESHOLD_USD:,.0f} threshold and this client "
            "cannot collect an approval.")
    return input_required(...)
\end{py}

Note what that does. It does not fail silently, and it does not send a shape the
client will choke on. It returns a tool execution error explaining exactly why
the operation cannot proceed, which the model can relay to the user, who can go
and use a client that supports elicitation.

\section{\texttt{requestState} is attacker-controlled}

Here is the part that deserves your full attention.

\texttt{requestState} leaves your server, passes through a client you do not
control, and comes back. It may come back modified. It may come back from a
different user. It may come back attached to a different tool call. It may come
back next week.

\begin{dangerbox}{Three attacks on a naive \texttt{requestState}}
\textbf{Tampering.} Server seals \texttt{\{"exposure": 9000000\}}. Client edits
it to \texttt{\{"exposure": 100\}} to slip under the approval threshold.

\textbf{Cross-user replay.} Alice triggers an approval and captures the state.
Bob presents it and inherits her half-completed privileged operation.

\textbf{Cross-request replay.} State minted for \texttt{approve\_refund} is
presented on \texttt{transfer\_funds}, and the server accepts an approval that
was never given for that operation.
\end{dangerbox}

The specification's requirement: if \texttt{requestState} influences
authorization, resource access, or business logic, servers \textbf{must} protect
its integrity and \textbf{must} reject state that fails verification. Integrity
protection may be omitted only when tampering can cause nothing worse than the
request failing.

Meridian's implementation closes all three, plus a window:

\begin{py}
class StateSealer:
    """Authenticated, expiring, principal-bound `requestState`.

    The format is `v1.<payload-b64>.<mac-b64>`, with the MAC over the payload
    bytes. HMAC-SHA256 gives integrity, which is what the specification
    requires. It does not give confidentiality: anyone holding the blob can
    read it. So put correlation data in here, never secrets.
    """

    def seal(self, data, *, principal=None, method=None, params=None,
             ttl_seconds=None, now=None) -> str:
        now = time.time() if now is None else now
        envelope = {
            "d": data,
            "exp": now + (ttl_seconds if ttl_seconds is not None
                          else self.ttl_seconds),
            "sub": principal,
            "req": self.request_digest(method, params or {}) if method else None,
        }
        payload = json.dumps(envelope, sort_keys=True,
                             separators=(",", ":")).encode()
        return (f"{self.VERSION}.{self._b64e(payload)}"
                f".{self._b64e(self._mac(payload))}")
\end{py}

Four bindings, each closing one door:

\begin{description}[leftmargin=1.4em, style=nextline, itemsep=3pt]
\item[The MAC] stops tampering. Verified in constant time, so a timing oracle
cannot walk the signature out of you.
\item[\texttt{sub}, the principal] stops cross-user replay.
\item[\texttt{req}, a request digest] stops cross-request replay. Note carefully
what it covers: the method, the name, and the arguments, but \emph{not}
\texttt{inputResponses}, because those change between the mint and the redeem.
\item[\texttt{exp}] bounds the window when the other three are not enough.
\end{description}

\begin{py}
if not hmac.compare_digest(mac, self._mac(payload)):
    raise StateTampered()
if envelope.get("exp", 0) < now:
    raise StateExpired()
bound_sub = envelope.get("sub")
if bound_sub is not None and bound_sub != principal:
    raise StateTampered("requestState was issued to a different principal")
bound_req = envelope.get("req")
if bound_req is not None:
    if not hmac.compare_digest(bound_req,
                               self.request_digest(method, params or {})):
        raise StateTampered("requestState was issued for a different request")
\end{py}

Each binding gets its own test, because a security control that is not tested is
a security control that will be refactored away:

\begin{py}
def test_principal_binding_blocks_cross_user_replay(self):
    token = self.sealer.seal({"txn": "T1"}, principal="alice")
    self.assertEqual(self.sealer.open(token, principal="alice"), {"txn": "T1"})
    with self.assertRaises(McpError):
        self.sealer.open(token, principal="bob")

def test_request_binding_blocks_cross_request_replay(self):
    params = {"name": "approve_refund", "arguments": {"amount": 10}}
    token = self.sealer.seal({"ok": True}, method="tools/call", params=params)
    other = {"name": "transfer_funds", "arguments": {"amount": 10}}
    with self.assertRaises(McpError):
        self.sealer.open(token, method="tools/call", params=other)
\end{py}

And an end-to-end test that checks both directions, because a check that rejects
everything is not a check:

\begin{py}
first = post({"name": "assess_account_risk", "arguments": args}, 1)["result"]
forged = first["requestState"][:-4] + "AAAA"
second = post({..., "requestState": forged}, 2)
self.assertIn("error", second, "a forged requestState was accepted")

# And the untampered state still works, so the check is not just
# rejecting everything.
third = post({..., "requestState": first["requestState"]}, 3)
self.assertIn("structuredContent", third["result"])
\end{py}

\begin{notebox}{None of this makes it single-use}
Expiry and binding bound the replay window and stop cross-user and
cross-request reuse. They do not guarantee at-most-once.

If an operation must happen at most once (a payment, a one-time redemption), the
server has to track consumption itself. There is no protocol feature that will
do it for you, and the specification says so explicitly.
\end{notebox}

Two more rules that fall out of the format. Sign with a secret shared across
your fleet, not a per-process one, or a retry that lands on another replica
cannot open state its sibling sealed. And put no secrets inside: HMAC gives
integrity, not confidentiality, and anyone holding the blob can read it.

\section{Elicitation}

The one input request you should actually build on. Two modes, and the choice
between them is a security decision.

\subsection{Form mode}

Structured data, collected through the client, validated against a schema. The
schema is restricted to flat objects of primitives, deliberately, so that any
client can render a form without implementing a JSON Schema editor.

Strings, numbers, integers, booleans, and enums (single and multi-select, with
or without display titles). No nested objects. No arrays of objects.

\begin{py}
"approval": elicit_form(
    f"Exposure of ${exposure:,.0f} on {account_id} exceeds the "
    f"${APPROVAL_THRESHOLD_USD:,.0f} threshold. Who is approving?",
    {
        "type": "object",
        "properties": {
            "approver": {
                "type": "string",
                "title": "Approver",
                "description": "Name or employee id of the approver",
                "minLength": 2,
            },
            "rationale": {"type": "string", "title": "Rationale",
                          "maxLength": 400},
        },
        "required": ["approver"],
    },
)
\end{py}

\begin{dangerbox}{Never use form mode for secrets}
The specification is a MUST NOT: no passwords, API keys, access tokens, or
payment credentials through form mode.

The reason is architectural. Form data passes through the client, is rendered in
a UI, and lands in logs. Contact details are a judgement call. Credentials are
not.

For anything sensitive, use URL mode.
\end{dangerbox}

\subsection{The three response actions}

\begin{table}[htbp]
\centering
\small
\begin{tabularx}{\textwidth}{@{}llX@{}}
\toprule
\thead{Action} & \thead{Means} & \thead{Server should} \\
\midrule
\texttt{accept}  & Submitted, with data & Process it \\
\texttt{decline} & Explicitly refused & Offer an alternative, or fail clearly \\
\texttt{cancel}  & Dismissed without choosing & Consider asking again later \\
\bottomrule
\end{tabularx}
\caption{\texttt{decline} and \texttt{cancel} are different signals and deserve
different handling. Collapsing them into ``no'' loses the distinction between
``I refuse'' and ``I pressed Escape''.}
\label{tab:elicit-actions}
\end{table}

\begin{py}
if answer.declined:
    return tool_error(
        f"Exposure of ${exposure:,.0f} was declined by the approver.")
if answer.cancelled:
    return tool_error("Approval dialog was dismissed. Nothing was scored.")
\end{py}

And one case that is neither: the user accepted but left a required field blank.
The specification's guidance is to \emph{ask again}, not to error:

\begin{py}
approver = (answer.content or {}).get("approver", "").strip()
if not approver:
    # Missing a value we need is not an error. Ask again.
    return _ask_for_approval(ctx, account_id, exposure)
\end{py}

\subsection{URL mode}

For anything that must not touch the client. The server returns a URL, the
client shows it, gets consent, and opens it somewhere it cannot read.

\begin{wire}
{
  "method": "elicitation/create",
  "params": {
    "mode": "url",
    "url": "https://mcp.meridian.example/ui/connect-bank",
    "message": "Connect your bank account to continue."
  }
}
\end{wire}

The response says only \texttt{\{"action": "accept"\}}, and that means ``the user
agreed to look'', not ``the interaction succeeded''. The client never learns the
outcome. The server finds out by other means and reports on the retry.

The client-side rules are strict, and they exist because a server-supplied URL
rendered in a trusted UI is a phishing vector:

\begin{itemize}[leftmargin=1.6em, itemsep=2pt]
\item Must not pre-fetch the URL or its metadata.
\item Must not open it without explicit consent.
\item Must show the full URL before consent.
\item Must open it where neither the client nor the model can inspect the
contents. On iOS that means \texttt{SFSafariViewController}, not
\texttt{WKWebView}.
\item Should highlight the domain, and warn on Punycode.
\end{itemize}

\begin{dangerbox}{The URL-mode phishing attack}
Alice, a legitimate user of a benign server, triggers a URL-mode elicitation
that starts a third-party OAuth flow. Her client shows the link. Instead of
clicking it, Alice sends it to Bob. Bob opens it, authorizes, and believes he
has connected his own account. The server receives the callback, assumes it
belongs to Alice's pending elicitation, and binds \emph{Bob's} third-party
tokens to \emph{Alice's} identity.

The fix is that the server must verify that the user who \emph{opens} the URL is
the user for whom it was minted. In practice: point the elicitation at your own
\texttt{/connect} route, check a session cookie against the \texttt{sub} claim
from your authorization server, and only then redirect onward.
\end{dangerbox}

\section{The cost}

Now the number, which is the reason this chapter exists in a performance book.

\bookfig[0.84]{xkcd-elicitation}{The transport cost of an elicitation is small.
The model turn it forces is not.}{xkcd-elicitation}

\begin{measurebox}{One elicitation, measured}
\begin{tabular}{@{}lrr@{}}
\toprule
\thead{} & \thead{Requests} & \thead{Latency} \\
\midrule
Tool call, no elicitation            & 1.0 & 34.3\,ms \\
Tool call with elicitation           & 2.0 & 68.8\,ms \\
\midrule
Extra transport                      &     & +34.5\,ms \\
Plus one model turn to read the answer &   & +340\,ms \\
\textbf{True cost}                   &     & \textbf{+375\,ms} \\
\bottomrule
\end{tabular}

\vspace{6pt}
\footnotesize
Measured at 12\,ms simulated per-call latency, the same run
Chapter~\ref{ch:primer} quoted. Repeated here because this is the chapter where
you decide whether to incur it.
\end{measurebox}

Plus the human, who is not in that table and who takes seconds.

\keyidea{Every MRTR round trip is at least one extra client round trip and
usually one extra model turn. The transport cost is around a tenth of the total.
Optimising the transport is missing the point; the answer is to not need the
round trip.}

\subsection{Four ways to not need it}

\textbf{Put it in the schema.} The best elicitation is the one that never
happens. If a tool always needs an approver for large exposures, make
\texttt{approver} an optional argument and let the model supply it when the user
already said it. Zero round trips when it works.

\textbf{Ask for everything at once.} \texttt{inputRequests} is a map. Three
questions in one interim result cost one round trip; three separate interim
results cost three.

\textbf{Ask early, not late.} An elicitation on step one of a five-step task
costs one model turn. On step four, it costs a model turn plus the risk that the
user has wandered off and the context has gone stale.

\textbf{Consider whether you need to ask at all.} Chapter~\ref{ch:tools} put it
as a question: should the tool schema have captured this? Often yes. The
elicitation is sometimes a symptom of a tool that was designed around an
existing API rather than around the task.

\begin{table}[htbp]
\centering
\small
\begin{tabularx}{\textwidth}{@{}lX@{}}
\toprule
\thead{Ask the...} & \thead{When} \\
\midrule
User, via elicitation
  & Only a human can supply it: authority, preference, consent, a credential \\
Model, via a better schema
  & It is derivable from the conversation the model is already holding \\
Nobody
  & The tool should have taken it as a parameter, and you are papering over a
    design gap \\
\bottomrule
\end{tabularx}
\caption{Elicitation versus prompting versus neither. The third row is more
common than people admit.}
\label{tab:ask-who}
\end{table}

\section{Why Sampling and Roots lost}

A design post-mortem, because the reasoning generalises.

\subsection{Sampling}

The idea: a server asks the client to run a completion on its behalf. The server
gets model access without an API key, the user's provider and quota are used,
and the host keeps control over what the model sees.

Genuinely elegant. It lost anyway, for three reasons.

\textbf{The control-flow inversion was the whole problem.} Sampling was the main
justification for server-initiated requests, and server-initiated requests are
what statelessness could not accommodate.

\textbf{Direct provider integration won on the ground.} It is 2026. A server
that wants a model call adds three lines and an API key. The indirection bought
less than it cost.

\textbf{Quality control is impossible through the indirection.} The server
cannot choose the model, cannot see the system prompt, and cannot verify what
actually ran. \texttt{modelPreferences} was a hint, and hints are not a contract.

\begin{legacybox}{Sampling}
Deprecated as of \texttt{2026-07-28}, eligible for removal no earlier than
2027-07-28. It still works, delivered through MRTR rather than as a server
request.

Maintaining a server that uses it: it functions unchanged during the window.

Migrating: call your provider's API directly. You will need your own key and
your own budget, and you will gain control over model choice and the ability to
test what you shipped.

Also deprecated alongside it: \texttt{includeContext: "thisServer"} and
\texttt{"allServers"}. Omit the field or use \texttt{"none"}.
\end{legacybox}

\subsection{Roots}

The idea: the client tells the server which directories it may work in.

It lost for a simpler reason. It was a capability that looked like a security
boundary and was not one. A server that respects roots is a server that was
going to behave anyway; a malicious one ignores the list. Meanwhile it added a
client-side concept, a notification, and a round trip.

\begin{legacybox}{Roots}
Deprecated as of \texttt{2026-07-28}, eligible for removal no earlier than
2027-07-28. \texttt{notifications/roots/list\_changed} was removed outright.

Migrate by passing directories and files as ordinary tool parameters, as
resource URIs, or through server configuration. All three are explicit, all
three are visible in an audit log, and none of them pretends to be a security
boundary.
\end{legacybox}

The generalisable lesson: \textbf{a protocol feature that depends on the other
party's goodwill is documentation, not enforcement.} Roots described an
intention. The host's consent gate enforces one.

\section{Meridian's approval flow, end to end}

Two round trips, and Meridian asserts exactly that:

\begin{py}
def test_approval_flow_completes_in_two_round_trips(self):
    host = wire_host(inputs={
        "approval": {"action": "accept",
                     "content": {"approver": "j.okonjo",
                                 "rationale": "board ok"}},
    })
    binding = host.bindings["risk"]
    before = binding.client.stats.requests

    result = host.call_tool("risk.assess_account_risk",
                            {"accountId": "ACC-1000", "exposureUsd": 9_000_000})

    self.assertFalse(result.get("isError"))
    self.assertIn("score", result["structuredContent"])
    self.assertEqual(binding.client.stats.requests - before, 2)
    self.assertEqual(binding.client.stats.mrtr_rounds, 1)
\end{py}

And crucially, the sub-threshold path costs one:

\begin{py}
def test_below_threshold_needs_no_round_trip(self):
    ...
    host.call_tool("risk.assess_account_risk",
                   {"accountId": "ACC-1000", "exposureUsd": 100_000})
    self.assertEqual(binding.client.stats.requests - before, 1)
\end{py}

That second test is the one that matters operationally. The expensive path
should be reachable only when it is genuinely needed, and a test is the only
thing that stops a refactor from making it unconditional.

The client-side loop is bounded, because a server can legitimately ask again and
a buggy one can ask forever:

\begin{py}
for round_no in range(MAX_MRTR_ROUNDS):
    ...
    if result.get("resultType") != jsonrpc.RESULT_INPUT_REQUIRED:
        return result
    self.stats.mrtr_rounds += 1
    requests = result.get("inputRequests") or {}
    pending_state = result.get("requestState")
    pending_responses = self._answer(requests) if requests else {}
    if requests and not pending_responses:
        # Nothing could be answered. Retrying would loop forever.
        raise errors.InvalidParams(
            "server requested input this client cannot provide")

raise errors.InternalError(
    f"{method} still asking for input after {MAX_MRTR_ROUNDS} rounds")
\end{py}

\section{What to remember}

\begin{itemize}[leftmargin=1.6em, itemsep=3pt]
\item Servers answer with \texttt{input\_required} instead of asking mid-call.
The original request finishes.
\item The retry uses a new id, repeats the arguments, and echoes
\texttt{requestState} unmodified.
\item \texttt{requestState} is attacker-controlled. Seal it with a MAC, bind it
to the principal and the request, and expire it. None of that makes it
single-use.
\item Sign with a fleet-wide secret, and put no secrets inside.
\item Never send an input request the client did not declare support for.
Degrade with a clear tool error.
\item Form mode for ordinary input, URL mode for anything sensitive. Verify that
the user who opens a URL is the user it was minted for.
\item \texttt{decline} and \texttt{cancel} are different. A blank required field
means ask again, not error.
\item One elicitation costs about 375\,ms, of which only 35 is transport. Ask
for everything at once, ask early, or design so you need not ask.
\item Sampling and Roots are deprecated. Both depended on the other party's
goodwill or on a stateful connection.
\end{itemize}

Next: extensions, and the Tasks extension, which is what you use when the work
takes longer than anyone wants to hold a connection open for.
